% \restoregeometry

% \newpage

\newpage

\thispagestyle{empty}

% \newgeometry{left=2.2cm,right=2.8cm,top=2cm,bottom=2cm}
% \pagenumbering{roman}
% \setcounter{page}{1}
% \thisfancypage{
% 	\setlength{\fboxsep}{0.5cm}
% 	\fbox}{}

\begin{tikzpicture}[remember picture, overlay]
\path ($(current page text area.north east)+(0cm, 0.5cm)$) coordinate (A)
($(current page text area.north west)+(0cm, 0.5cm)$) coordinate (B)
(current page text area.south west) coordinate (C)
(current page text area.south east) coordinate (D);
\draw[black!40!black, line width=3pt]
(A)--(B)--(C)--(D)--cycle;
\draw[black!10!black, line width=0.5pt]
($(A)+(-0.1cm, -0.1cm)$)--($(B)+(0.1cm, -0.1cm)$)--($(C)+(0.1cm, 0.1cm)$)--($(D)+(-0.1cm, 0.1cm)$)--cycle;
\end{tikzpicture}

\fontsize{13pt}{16pt}\selectfont
\bigbreak
\begin{center}
{\bfseries NHẬN XÉT CỦA GIẢNG VIÊN HƯỚNG DẪN}

\end{center}
\bigbreak

\fontsize{12pt}{14pt}\selectfont
\begin{enumerate}
\item [{\bfseries 1.}]{\bfseries Mục đích và nội dung của đồ án}
\begin{enumerate}
\item \item \item \end{enumerate}
\item [{\bfseries 2.}] {\bfseries Kết quả đạt được}
\begin{enumerate}
\item \item \item \end{enumerate}
\item [{\bfseries 3.}]{\bfseries Ý thức làm việc của sinh viên}
\begin{enumerate}
\item \item \item \end{enumerate}
\end{enumerate}
\hspace{0.4\textwidth}
\begin{minipage}{0.5\textwidth}
\noindent\begin{center}
% \textit{Hà Nội, ngày 06 tháng 02 năm 2022} \\
\textit{Hà Nội, ngày \dots tháng \dots năm \dots} \\

Giảng viên hướng dẫn \\
\textit{(Ký và ghi rõ họ tên)} \\
\vspace{2cm}
% \textbf{PGS.TS. Nguyễn Thị Thu Thủy}
\textbf{TS. Trần Anh Tú}

\end{center}
\end{minipage}
% \restoregeometry

\newpage
